\renewcommand{\X}{|[fill=black,   text=white]|+}
\renewcommand{\O}{|[fill=white,   text=black]|O}
\renewcommand{\d}{|[fill=gray!50 ,text=ocre] |.}
\nsection{Yin Yang}


The puzzle {\it Yin Yang}~:
\wwwurl{https://www.puzzle-yin-yang.com/}

\noindent consists of a $2D$ rectangular grid of cells - black ('+'), white ('O') or
unknown ('.').
\begin{center}
\begin{tikzpicture}[every node/.style={draw=black,anchor=base,text depth=.5ex,text height=2ex,text width=1em,outer sep=0pt,align=center,inner sep=0pt}] \matrix [matrix of nodes,draw=white,nodes in empty cells] {
\X&\d&\X&\d&\X\\
\d&\d&\O&\d&\X\\
\O&\X&\X&\X&\d\\
\O&\d&\d&\O&\X\\
\X&\d&\X&\d&\d\\
};
\end{tikzpicture}
\end{center}
The aim is to change all the unknowns to either black or white
such that:

\begin{enumerate}
\item All black cells should be $4$-connected together (up, down,
left, right) into one single region.
\item All white cells should be $4$-connected together (up, down,
left, right) into one single region.
\item $2\times 2$ areas of the same colour are not allowed.
\begin{center}
\begin{tikzpicture}[every node/.style={draw=black,anchor=base,text depth=.5ex,text height=2ex,text width=1em,outer sep=0pt,align=center,inner sep=0pt}] \matrix [matrix of nodes,draw=white,nodes in empty cells] {
\X&\X\\
\X&\X\\
};
\end{tikzpicture}
\hspace*{0.5in}
\begin{tikzpicture}[every node/.style={draw=black,anchor=base,text depth=.5ex,text height=2ex,text width=1em,outer sep=0pt,align=center,inner sep=0pt}] \matrix [matrix of nodes,draw=white,nodes in empty cells] {
\O&\O\\
\O&\O\\
};
\end{tikzpicture}
\end{center}
\end{enumerate}

\noindent From these somewhat simple rules, another constraint emerges, that no
$2\times 2$ area can have diagonally opposed colours. Therefore, 
the following $2\times 2$ areas are not allowed either~:
\begin{center}
\begin{tikzpicture}[every node/.style={draw=black,anchor=base,text depth=.5ex,text height=2ex,text width=1em,outer sep=0pt,align=center,inner sep=0pt}] \matrix [matrix of nodes,draw=white,nodes in empty cells] {
\X&\O\\
\O&\X\\
};
\end{tikzpicture}
\hspace*{0.5in}
\begin{tikzpicture}[every node/.style={draw=black,anchor=base,text depth=.5ex,text height=2ex,text width=1em,outer sep=0pt,align=center,inner sep=0pt}] \matrix [matrix of nodes,draw=white,nodes in empty cells] {
\O&\X\\
\X&\O\\
};
\end{tikzpicture}
\end{center}

One approach to solving such a puzzle is via a brute-force algorithm, exploring all possible
(valid) combinations of black and white cells. Doing this can lead to the correct solution
of the grid above~:
\begin{center}
\begin{tikzpicture}[every node/.style={draw=black,anchor=base,text depth=.5ex,text height=2ex,text width=1em,outer sep=0pt,align=center,inner sep=0pt}] \matrix [matrix of nodes,draw=white,nodes in empty cells] {
\X&\X&\X&\X&\X\\
\O&\O&\O&\O&\X\\
\O&\X&\X&\X&\X\\
\O&\O&\O&\O&\X\\
\X&\X&\X&\X&\X\\
};
\end{tikzpicture}
\end{center}


\begin{exercise}
Write the functions specified in \verb^yy.h^ that allows grids to
be read from file, converted to a string, and the number
of solutions available computed.

\noindent
The brute-force algorithm for searching all moves for a
solution goes like this~:
\begin{enumerate}
\item You will use a large array (list) of grids.
\item Put the initial grid into the front of this list, \verb^f=0^.
\item Consider the grid at the {\bf front} of the list (index \verb$f$).
\item For this (parent) grid, create the two `child' grids which can be made 
by filling the {\bf next} unknown cell with either black or white.
\begin{itemize}
\item If this grid is valid (see $2\times 2$ rules above), add it to the end of the list.
\item If this is the `final' grid (i.e. all unknown cells have been filled in and there is
exactly one contiguous black region and exactly one contiguous white region), then add
one to the number of possible (valid) solutions.
\end{itemize}
\item Add one to $f$. If there are more grids in the list, go to step $3$.
\item If not, return the number of possible solutions found.
\end{enumerate}

\noindent
Your program~:
\begin{itemize}
\item {\bf Should} use the algorithm detailed above (which is similar to a queue and therefore a breadth-first search). Do not use the other algorithms possible; the quality of your coding is being assessed against others taking the same approach, and if you do something different it won't get any marks.
\item {\bf Should} run silently (except in the case of an \verb^EXIT_FAILURE^, maybe due to an \verb^alloc()^ error).
\item {\bf Should} implement the function \verb^test()^ to perform any assertion testing, (including tests of your helper functions). If you create new grid files (other than the ones I provde) make sure these are submitted as part of your \verb^.zip^ file.
\item Should {\bf NOT} use linked lists.
The list of grids can be a fixed (large) size (I used $250,000$)
\item Should {\bf NOT} print anything out to screen after successfully
completing the search.
\end{itemize}


\subsection*{Extension}

Basic assignment = {\Large $90\%$}.
Extension = {\Large $10\%$}.

\noindent
{\bf If} you get your program working, there's a small extension available.
Recode your program so that it executes as quickly as possible in files
called \verb^ext.c^ and \verb^ext_mydefs.h^.
For this I will compute the time taken for your code to solve
a series of complex grids using \verb^/usr/bin/time^.
No algorithm is specified for this and no coding rules need to be followed,  
except that it must compile using the flags I've specified (without warnings).
\end{exercise}
